\chapter{BX-64: Application Vectors with V and Zve64d}
\label{ch:bx64}

\section{The architecture}

BX-64 widens XLEN to 64. The base is RV64I. Chapter~\ref{ch:rv64} is that integer file. The vector unit is Zve64d or the full application extension V. ELEN is 64. VLEN is at least 128, and 256 is the recommended elaboration when the block RAM budget allows a file of thirty-two registers times 256 bits, which is one kilobyte.

\begin{spec}
Legal SEW = 8, 16, 32, and 64.\\
FP32 and FP64, including fused multiply-add.\\
If the claim is V: every instruction V requires, VLEN at least 128,\\
and the V bit set in \texttt{misa}.\\
If the claim is Zve64d: do not set the V bit.
\end{spec}

The second sentence is the one that keeps the ladder honest. A machine that implements double-precision vector arithmetic and not yet the whole of V is a useful machine. It is not a V machine. Chapters~\ref{ch:bx128} and~\ref{ch:bx512} keep this element width and pin the register length. An integer register of 128 bits is Chapter~\ref{ch:rv128}.

\section{What arrives with sixty-four-bit elements}

SEW=64 is the element width a double-precision value needs, and the element width a widened thirty-two-bit product needs if you choose to widen that far. Both appear in geometry and in audio rate conversion. The lane's integer ALU grows a sixty-four-bit path, which on a sixteen-bit or thirty-two-bit DSP primitive is a short sequence of partial adds. That sequence is internal. The architectural latency is a published integer.

Masks, slides, and reductions at SEW=64 are required. Indexed accesses are required. \texttt{vstart} becomes architecturally visible if the machine has a page table and a vector load can fault on a later page than it started on. The simple machines avoided this by checking the whole access first. A sixty-four-bit hart with virtual memory cannot always do that cheaply. Implement \texttt{vstart}, test a fault on element 0 and a fault on a middle element, and document that restart is precise.

\section{Opcode map for Zve64d / V (BX-64)}

All opcodes of Zve32f remain legal. The additional element widths and FP64 operations use the same rigid two-line matrix. SEW may be 8, 16, 32 or 64. Every major class has an explicit entry diagram below or inherits the diagram from the lower profile.

\subsection{Configuration}

Identical to the lower profiles (see Chapter~\ref{ch:bx8}).

\subsection{Unit-stride load at SEW=64 --- \texttt{vle64.v}}

\begin{center}
\ttfamily\footnotesize
\setlength{\tabcolsep}{0.3em}
\begin{tabular}{*{16}{w{c}{1.05em}}}
31 & 30 & 29 & 28 & 27 & 26 & 25 & 24 & 23 & 22 & 21 & 20 & 19 & 18 & 17 & 16 \\
\hline
\multicolumn{3}{|c|}{N} &
\multicolumn{1}{c|}{E} &
\multicolumn{2}{c|}{P} &
\multicolumn{1}{c|}{M} &
\multicolumn{5}{c|}{L} &
\multicolumn{4}{c|}{S} \\
\hline
\end{tabular}
\vspace{0.35ex}
\begin{tabular}{*{16}{w{c}{1.05em}}}
15 & 14 & 13 & 12 & 11 & 10 & 9 & 8 & 7 & 6 & 5 & 4 & 3 & 2 & 1 & 0 \\
\hline
\multicolumn{3}{|c|}{W} &
\multicolumn{5}{c|}{D} &
\multicolumn{7}{c|}{O} \\
\hline
\end{tabular}
\end{center}

\vspace{1.4ex}
\begin{center}
\begin{tabular}{@{}ll@{}}
N & = nf \\
E & = mew \\
P & = mop (00 = unit-stride) \\
M & = vm \\
L & = lumop (00000) \\
S & = rs1 (base) \\
W & = width (111 = 64-bit) \\
D & = vd \\
O & = 0000111 \\
\end{tabular}
\end{center}

Stride and indexed forms keep the same layout; only the mop and width fields change (width = 111 for EEW=64).

\subsection{Integer / FP arithmetic (OP-V) at any SEW including 64}

\begin{center}
\ttfamily\footnotesize
\setlength{\tabcolsep}{0.3em}
\begin{tabular}{*{16}{w{c}{1.05em}}}
31 & 30 & 29 & 28 & 27 & 26 & 25 & 24 & 23 & 22 & 21 & 20 & 19 & 18 & 17 & 16 \\
\hline
\multicolumn{6}{|c|}{F} &
\multicolumn{1}{c|}{M} &
\multicolumn{5}{c|}{S} &
\multicolumn{4}{c|}{T} \\
\hline
\end{tabular}
\vspace{0.35ex}
\begin{tabular}{*{16}{w{c}{1.05em}}}
15 & 14 & 13 & 12 & 11 & 10 & 9 & 8 & 7 & 6 & 5 & 4 & 3 & 2 & 1 & 0 \\
\hline
\multicolumn{1}{|c|}{T} &
\multicolumn{3}{c|}{C} &
\multicolumn{5}{c|}{D} &
\multicolumn{7}{c|}{O} \\
\hline
\end{tabular}
\end{center}

\vspace{1.4ex}
\begin{center}
\begin{tabular}{@{}ll@{}}
F & = funct6 \\
M & = vm \\
S & = vs2 \\
T & = vs1 / rs1 / fs1 \\
C & = funct3 (OPIVV / OPIVX / OPFVV / OPFVF \ldots) \\
D & = vd \\
O & = 1010111 \\
\end{tabular}
\end{center}

All integer instructions of Zve32f are legal at SEW=64. All FP32 instructions remain legal; their FP64 counterparts use the same funct6 values with SEW=64.

\begin{center}
\begin{tabular}{@{}llc@{}}
\toprule
Class & examples & notes \\
\midrule
Integer & \texttt{vadd}, \texttt{vwmul}, \texttt{vredsum}, \texttt{vrgather} & full set \\
FP32   & \texttt{vfadd}, \texttt{vfmacc}, \texttt{vfcvt} & from BX-32 \\
FP64   & \texttt{vfadd}, \texttt{vfmacc}, \texttt{vfcvt} at SEW=64 & required for Zve64d/V \\
\bottomrule
\end{tabular}
\end{center}

If the claim is the full application extension V, every additional instruction required by the V specification must be present and the V bit set in \texttt{misa}. A Zve64d implementation omits the V bit and may omit instructions that V adds beyond Zve64d.

\subsection{Mask and move}

Full mask set plus the FP move variants for both single and double.

Xsvm must be left off on any image that claims V.

\section{Assembler kernels that must run}

In addition to every kernel of the lower profiles, the following must pass on BX-64. The COP3 framing of Chapter~\ref{ch:bx32} is the host pattern at this width too. The address is the low 64 bits of the GPR. There is no separate harness directory in this edition.

\begin{table}[htbp]
\centering
\caption{Kernel analogues on this machine.}
\small
\begin{tabular}{@{}llll@{}}
\toprule
Kernel & RVV & Intel & ARM \\
\midrule
sew64 & \texttt{vse64.v} & \texttt{vmovupd} & \texttt{st1} \\
fma64 & \texttt{vfmacc.vv} & \texttt{vfmadd231pd} & \texttt{fmla} \\
gemm64 & \texttt{vfmacc.vf} & \texttt{vmovddup} & \texttt{ld1r} \\
index64 & \texttt{vluxei64.v} & \texttt{vpgatherqq} & \texttt{ldr} \\
\bottomrule
\end{tabular}
\end{table}

The comparison table under each kernel is one row per instruction: RVV, then Intel, then ARM. Chapter~\ref{ch:bx8} says which encodings those bytes are.

\subsection{sew64 --- double-word fill}
The element is 64 bits. VLEN 128 holds two of them.

\begin{table}[htbp]
\centering
\caption{Kernel comparison: sew64}
\begingroup
\scriptsize
\setlength{\tabcolsep}{4pt}
\noindent \texttt{a0 = n, a1 = fill, a2 = dst}\par
\vspace{0.6ex}
\begin{tabular}{@{}l|l|l|p{7.2cm}@{}}
\toprule
ISA & Address & Hex & Assembly \\
\midrule
RVV & \texttt{0000} & \texttt{02050063} & \texttt{sew64: beqz a0, end} \\
RVV & \texttt{0004} & \texttt{003572D7} & \texttt{loop: vsetvli t0, a0, e64, m1, tu, mu} \\
RVV & \texttt{0008} & \texttt{5E05C0D7} & \texttt{vmv.v.x v1, a1} \\
RVV & \texttt{000C} & \texttt{020670A7} & \texttt{vse64.v v1, (a2)} \\
RVV & \texttt{0010} & \texttt{00329313} & \texttt{slli t1, t0, 3} \\
RVV & \texttt{0014} & \texttt{00660633} & \texttt{add a2, a2, t1} \\
RVV & \texttt{0018} & \texttt{40550533} & \texttt{sub a0, a0, t0} \\
RVV & \texttt{001C} & \texttt{FE0514E3} & \texttt{bnez a0, loop} \\
RVV & \texttt{0020} & \texttt{00008067} & \texttt{end: ret} \\
\midrule
Intel & \texttt{0000} & \texttt{C4E1F96EC6} & \texttt{sew64: vmovq xmm0, rsi} \\
Intel & \texttt{0005} & \texttt{C4E27959C0} & \texttt{vpbroadcastq xmm0, xmm0} \\
Intel & \texttt{000A} & \texttt{C5F91107} & \texttt{vmovupd [rdi], xmm0} \\
\midrule
ARM & \texttt{0000} & \texttt{4E080C20} & \texttt{sew64: dup v0.2d, x1} \\
ARM & \texttt{0004} & \texttt{4C007C40} & \texttt{st1 \{v0.2d\}, [x2]} \\
\bottomrule
\end{tabular}
\endgroup
\end{table}
\noindent\path{kernels/bx64/sew64.sew64.lmul1.tu.mu.vlen128.xlen64.le.bin}\par


\subsection{fma64 --- double-precision fused multiply-add}
\begin{table}[htbp]
\centering
\caption{Kernel comparison: fma64}
\begingroup
\scriptsize
\setlength{\tabcolsep}{4pt}
\noindent \texttt{a0 = n, a1 = a, a2 = b, a3 = c. c = c + a*b}\par
\vspace{0.6ex}
\begin{tabular}{@{}l|l|l|p{7.2cm}@{}}
\toprule
ISA & Address & Hex & Assembly \\
\midrule
RVV & \texttt{0000} & \texttt{02050A63} & \texttt{fma64: beqz a0, end} \\
RVV & \texttt{0004} & \texttt{003572D7} & \texttt{loop: vsetvli t0, a0, e64, m1, tu, mu} \\
RVV & \texttt{0008} & \texttt{0205F087} & \texttt{vle64.v v1, (a1)} \\
RVV & \texttt{000C} & \texttt{02067107} & \texttt{vle64.v v2, (a2)} \\
RVV & \texttt{0010} & \texttt{0206F187} & \texttt{vle64.v v3, (a3)} \\
RVV & \texttt{0014} & \texttt{B22091D7} & \texttt{vfmacc.vv v3, v1, v2} \\
RVV & \texttt{0018} & \texttt{0206F1A7} & \texttt{vse64.v v3, (a3)} \\
RVV & \texttt{001C} & \texttt{00329313} & \texttt{slli t1, t0, 3} \\
RVV & \texttt{0020} & \texttt{006585B3} & \texttt{add a1, a1, t1} \\
RVV & \texttt{0024} & \texttt{00660633} & \texttt{add a2, a2, t1} \\
RVV & \texttt{0028} & \texttt{006686B3} & \texttt{add a3, a3, t1} \\
RVV & \texttt{002C} & \texttt{40550533} & \texttt{sub a0, a0, t0} \\
RVV & \texttt{0030} & \texttt{FC051AE3} & \texttt{bnez a0, loop} \\
RVV & \texttt{0034} & \texttt{00008067} & \texttt{end: ret} \\
\midrule
Intel & \texttt{0000} & \texttt{C5FD1006} & \texttt{fma64: vmovupd ymm0, [rsi]} \\
Intel & \texttt{0004} & \texttt{C5FD100A} & \texttt{vmovupd ymm1, [rdx]} \\
Intel & \texttt{0008} & \texttt{C5FD1011} & \texttt{vmovupd ymm2, [rcx]} \\
Intel & \texttt{000C} & \texttt{C4E2FDB8D1} & \texttt{vfmadd231pd ymm2, ymm0, ymm1} \\
Intel & \texttt{0011} & \texttt{C5FD1111} & \texttt{vmovupd [rcx], ymm2} \\
\midrule
ARM & \texttt{0000} & \texttt{4C407C20} & \texttt{fma64: ld1 \{v0.2d\}, [x1]} \\
ARM & \texttt{0004} & \texttt{4C407C41} & \texttt{ld1 \{v1.2d\}, [x2]} \\
ARM & \texttt{0008} & \texttt{4C407C62} & \texttt{ld1 \{v2.2d\}, [x3]} \\
ARM & \texttt{000C} & \texttt{4E61CC02} & \texttt{fmla v2.2d, v0.2d, v1.2d} \\
ARM & \texttt{0010} & \texttt{4C007C62} & \texttt{st1 \{v2.2d\}, [x3]} \\
\bottomrule
\end{tabular}
\endgroup
\end{table}
\noindent\path{kernels/bx64/fma64.sew64.lmul1.tu.mu.vlen128.xlen64.le.bin}\par


\subsection{gemm64 --- outer product of a 2-wide panel}

The BX-32 catalogue in Section~\ref{sec:blas} is the same list at SEW 64, with \texttt{fma64}, \texttt{index64}, and \texttt{gemm64}. VLEN 128 holds two doubles, so the acceptance tile is 2 by 4: one vector of A and four broadcast scalars from B. \texttt{fld} is LOAD-FP with funct3 011.

The GEMM note's AVX2 tile holds four doubles in a \texttt{ymm}. Its AVX-512 tile is 24 by 8, eight doubles in each \texttt{zmm}. The Intel column here is two doubles in an \texttt{xmm}. \texttt{vmovddup} writes both lanes from one memory double.

\begin{table}[htbp]
\centering
\caption{Kernel comparison: gemm64}
\begingroup
\scriptsize
\setlength{\tabcolsep}{4pt}
\noindent \texttt{a0 = K, a1 = A panel, a2 = B panel, a3 = C. 2-wide outer product.}\par
\vspace{0.6ex}
\begin{tabular}{@{}l|l|l|p{7.2cm}@{}}
\toprule
ISA & Address & Hex & Assembly \\
\midrule
RVV & \texttt{0000} & \texttt{06050A63} & \texttt{gemm64: beqz a0, end} \\
RVV & \texttt{0004} & \texttt{00307057} & \texttt{vsetvli zero, zero, e64, m1, tu, mu} \\
RVV & \texttt{0008} & \texttt{0206F407} & \texttt{vle64.v v8, (a3)} \\
RVV & \texttt{000C} & \texttt{01068293} & \texttt{addi t0, a3, 16} \\
RVV & \texttt{0010} & \texttt{0202F487} & \texttt{vle64.v v9, (t0)} \\
RVV & \texttt{0014} & \texttt{02068293} & \texttt{addi t0, a3, 32} \\
RVV & \texttt{0018} & \texttt{0202F507} & \texttt{vle64.v v10, (t0)} \\
RVV & \texttt{001C} & \texttt{03068293} & \texttt{addi t0, a3, 48} \\
RVV & \texttt{0020} & \texttt{0202F587} & \texttt{vle64.v v11, (t0)} \\
RVV & \texttt{0024} & \texttt{0205F087} & \texttt{loop: vle64.v v1, (a1)} \\
RVV & \texttt{0028} & \texttt{00063007} & \texttt{fld ft0, 0(a2)} \\
RVV & \texttt{002C} & \texttt{B2105457} & \texttt{vfmacc.vf v8, ft0, v1} \\
RVV & \texttt{0030} & \texttt{00863007} & \texttt{fld ft0, 8(a2)} \\
RVV & \texttt{0034} & \texttt{B21054D7} & \texttt{vfmacc.vf v9, ft0, v1} \\
RVV & \texttt{0038} & \texttt{01063007} & \texttt{fld ft0, 16(a2)} \\
RVV & \texttt{003C} & \texttt{B2105557} & \texttt{vfmacc.vf v10, ft0, v1} \\
RVV & \texttt{0040} & \texttt{01863007} & \texttt{fld ft0, 24(a2)} \\
RVV & \texttt{0044} & \texttt{B21055D7} & \texttt{vfmacc.vf v11, ft0, v1} \\
RVV & \texttt{0048} & \texttt{01058593} & \texttt{addi a1, a1, 16} \\
RVV & \texttt{004C} & \texttt{02060613} & \texttt{addi a2, a2, 32} \\
RVV & \texttt{0050} & \texttt{FFF50513} & \texttt{addi a0, a0, -1} \\
RVV & \texttt{0054} & \texttt{FC0518E3} & \texttt{bnez a0, loop} \\
RVV & \texttt{0058} & \texttt{0206F427} & \texttt{vse64.v v8, (a3)} \\
RVV & \texttt{005C} & \texttt{01068293} & \texttt{addi t0, a3, 16} \\
RVV & \texttt{0060} & \texttt{0202F4A7} & \texttt{vse64.v v9, (t0)} \\
RVV & \texttt{0064} & \texttt{02068293} & \texttt{addi t0, a3, 32} \\
RVV & \texttt{0068} & \texttt{0202F527} & \texttt{vse64.v v10, (t0)} \\
RVV & \texttt{006C} & \texttt{03068293} & \texttt{addi t0, a3, 48} \\
RVV & \texttt{0070} & \texttt{0202F5A7} & \texttt{vse64.v v11, (t0)} \\
RVV & \texttt{0074} & \texttt{00008067} & \texttt{end: ret} \\
\midrule
Intel & \texttt{0000} & \texttt{C5F91006} & \texttt{gemm64: vmovupd xmm0, [rsi]} \\
Intel & \texttt{0004} & \texttt{C5FB120A} & \texttt{vmovddup xmm1, qword ptr [rdx]} \\
Intel & \texttt{0008} & \texttt{C4E2F9B8E1} & \texttt{vfmadd231pd xmm4, xmm0, xmm1} \\
Intel & \texttt{000D} & \texttt{C5FB124A08} & \texttt{vmovddup xmm1, qword ptr [rdx+8]} \\
Intel & \texttt{0012} & \texttt{C4E2F9B8E9} & \texttt{vfmadd231pd xmm5, xmm0, xmm1} \\
Intel & \texttt{0017} & \texttt{C5FB124A10} & \texttt{vmovddup xmm1, qword ptr [rdx+16]} \\
Intel & \texttt{001C} & \texttt{C4E2F9B8F1} & \texttt{vfmadd231pd xmm6, xmm0, xmm1} \\
Intel & \texttt{0021} & \texttt{C5FB124A18} & \texttt{vmovddup xmm1, qword ptr [rdx+24]} \\
Intel & \texttt{0026} & \texttt{C4E2F9B8F9} & \texttt{vfmadd231pd xmm7, xmm0, xmm1} \\
Intel & \texttt{002B} & \texttt{4883C610} & \texttt{add rsi, 16} \\
Intel & \texttt{002F} & \texttt{4883C220} & \texttt{add rdx, 32} \\
Intel & \texttt{0033} & \texttt{4883E901} & \texttt{sub rcx, 1} \\
Intel & \texttt{0037} & \texttt{75C7} & \texttt{jnz gemm64} \\
\midrule
ARM & \texttt{0000} & \texttt{4CDF7C20} & \texttt{gemm64: ld1 \{v0.2d\}, [x1], \#16} \\
ARM & \texttt{0004} & \texttt{4DDFCC41} & \texttt{ld1r \{v1.2d\}, [x2], \#8} \\
ARM & \texttt{0008} & \texttt{4E61CC04} & \texttt{fmla v4.2d, v0.2d, v1.2d} \\
ARM & \texttt{000C} & \texttt{4DDFCC41} & \texttt{ld1r \{v1.2d\}, [x2], \#8} \\
ARM & \texttt{0010} & \texttt{4E61CC05} & \texttt{fmla v5.2d, v0.2d, v1.2d} \\
ARM & \texttt{0014} & \texttt{4DDFCC41} & \texttt{ld1r \{v1.2d\}, [x2], \#8} \\
ARM & \texttt{0018} & \texttt{4E61CC06} & \texttt{fmla v6.2d, v0.2d, v1.2d} \\
ARM & \texttt{001C} & \texttt{4DDFCC41} & \texttt{ld1r \{v1.2d\}, [x2], \#8} \\
ARM & \texttt{0020} & \texttt{4E61CC07} & \texttt{fmla v7.2d, v0.2d, v1.2d} \\
ARM & \texttt{0024} & \texttt{F1000400} & \texttt{subs x0, x0, \#1} \\
ARM & \texttt{0028} & \texttt{54FFFEC1} & \texttt{b.ne gemm64} \\
\bottomrule
\end{tabular}
\endgroup
\end{table}
\noindent\path{kernels/bx64/gemm64.sew64.lmul1.tu.mu.vlen128.xlen64.le.bin}\par


\subsection{index64 --- indexed load at SEW=64}
The index EEW and the data EEW are both 64.

\begin{table}[htbp]
\centering
\caption{Kernel comparison: index64}
\begingroup
\scriptsize
\setlength{\tabcolsep}{4pt}
\noindent \texttt{a0 = n, a1 = base, a2 = index vector, a3 = dst}\par
\vspace{0.6ex}
\begin{tabular}{@{}l|l|l|p{7.2cm}@{}}
\toprule
ISA & Address & Hex & Assembly \\
\midrule
RVV & \texttt{0000} & \texttt{02050463} & \texttt{index64: beqz a0, end} \\
RVV & \texttt{0004} & \texttt{003572D7} & \texttt{loop: vsetvli t0, a0, e64, m1, tu, mu} \\
RVV & \texttt{0008} & \texttt{02067107} & \texttt{vle64.v v2, (a2)} \\
RVV & \texttt{000C} & \texttt{0625F087} & \texttt{vluxei64.v v1, (a1), v2} \\
RVV & \texttt{0010} & \texttt{0206F0A7} & \texttt{vse64.v v1, (a3)} \\
RVV & \texttt{0014} & \texttt{00329313} & \texttt{slli t1, t0, 3} \\
RVV & \texttt{0018} & \texttt{00660633} & \texttt{add a2, a2, t1} \\
RVV & \texttt{001C} & \texttt{006686B3} & \texttt{add a3, a3, t1} \\
RVV & \texttt{0020} & \texttt{40550533} & \texttt{sub a0, a0, t0} \\
RVV & \texttt{0024} & \texttt{FE0510E3} & \texttt{bnez a0, loop} \\
RVV & \texttt{0028} & \texttt{00008067} & \texttt{end: ret} \\
\midrule
Intel & \texttt{0000} & \texttt{C4E2ED91040E} & \texttt{index64: vpgatherqq ymm0, [rsi + ymm1*1], ymm2} \\
\midrule
ARM & \texttt{0000} & \texttt{FC637820} & \texttt{index64: ldr d0, [x1, x3, lsl \#3]} \\
ARM & \texttt{0004} & \texttt{FD000040} & \texttt{str d0, [x2]} \\
ARM & \texttt{0008} & \texttt{91002063} & \texttt{add x3, x3, \#8} \\
ARM & \texttt{000C} & \texttt{91002042} & \texttt{add x2, x2, \#8} \\
ARM & \texttt{0010} & \texttt{F1000400} & \texttt{subs x0, x0, \#1} \\
ARM & \texttt{0014} & \texttt{54FFFF61} & \texttt{b.ne index64} \\
\bottomrule
\end{tabular}
\endgroup
\end{table}
\noindent\path{kernels/bx64/index64.sew64.lmul1.tu.mu.vlen128.xlen64.le.bin}\par


\section{The same bodies at SEW 64}

The names in Section~\ref{sec:blas} with a \texttt{d} prefix use these images. \texttt{ddot} is \texttt{dot64}. \texttt{dnrm2} is \texttt{nrm264}. \texttt{dasum} is \texttt{asum64}. \texttt{drotg} and \texttt{drotmg} stay scalar. \texttt{idamax} is one strip of \texttt{iamax64}. The product is \texttt{gemm64}. The Intel column of that product is two doubles, via \texttt{vmovddup}. A \texttt{ymm} holds four doubles. A 24 by 8 tile is \texttt{zmm}, and that register file is Chapter~\ref{ch:bx512}.

Reassemble \texttt{bfly32}, \texttt{wmul32}, and \texttt{cmul32} with \texttt{e64} when the samples are doubles. FFTW's default transform, NumPy's \texttt{fft}, and a double-precision channelizer are that butterfly. GNU Radio's fast path stays at SEW 32.

\begin{table}[htbp]
\centering
\caption{Programs that run the BX-64 kernels.}
\begingroup
\scriptsize
\setlength{\tabcolsep}{3pt}
\begin{tabular}{@{}llp{7.0cm}@{}}
\toprule
Kernel & Work & Where it ships \\
\midrule
sew64 & doubleword fill & A zeroed double buffer. \\
fma64 & element $c + ab$ & A double polynomial step. The inner product of a compensated sum, written as a vector. \\
gemm64 & 2-wide rank-1 & \texttt{dgemm} in OpenBLAS, BLIS, and LAPACK's callers. The GEMM note's double update, here two doubles wide. \\
index64 & gather of doubles & Sparse matrix-vector with double values. A gather of double-precision samples. \\
scal64 & scale & \texttt{dscal}. Column scaling in LAPACK. \\
copy64 & copy & \texttt{dcopy}. Packing a double panel. \\
axpy64 & scaled add & \texttt{daxpy}. CG in double precision. \\
dot64 & inner product & \texttt{ddot}. A double correlation. SoXR's high-quality FIR. \\
nrm264 & Euclidean length & \texttt{dnrm2}. A residual norm in a double least-squares solve. \\
asum64 & L1 norm & \texttt{dasum}. \\
swap64 & exchange & \texttt{dswap}. Partial pivoting. \\
rot64 & Givens rotation & \texttt{drot}. Double QR. \\
gemv64 & matrix-vector & \texttt{dgemv}. The GEMM note's \texttt{cblas\_dgemv}, two doubles tall. \\
ger64 & rank-1 & \texttt{dger}. A double covariance update. \\
iamax64 & index of max abs & \texttt{idamax}. The pivot of a double LU. One strip. \\
\bottomrule
\end{tabular}
\endgroup
\end{table}

\begin{table}[htbp]
\centering
\caption{Kernel comparison: scal64}
\begingroup
\scriptsize
\setlength{\tabcolsep}{4pt}
\noindent \texttt{a0 = n, a1 = x, a2 = alpha. x *= alpha}\par
\vspace{0.6ex}
\begin{tabular}{@{}l|l|l|p{7.2cm}@{}}
\toprule
ISA & Address & Hex & Assembly \\
\midrule
RVV & \texttt{0000} & \texttt{00063007} & \texttt{scal64: fld ft0, 0(a2)} \\
RVV & \texttt{0004} & \texttt{02050263} & \texttt{beqz a0, end} \\
RVV & \texttt{0008} & \texttt{003572D7} & \texttt{loop: vsetvli t0, a0, e64, m1, tu, mu} \\
RVV & \texttt{000C} & \texttt{0205F087} & \texttt{vle64.v v1, (a1)} \\
RVV & \texttt{0010} & \texttt{921050D7} & \texttt{vfmul.vf v1, v1, ft0} \\
RVV & \texttt{0014} & \texttt{0205F0A7} & \texttt{vse64.v v1, (a1)} \\
RVV & \texttt{0018} & \texttt{00329313} & \texttt{slli t1, t0, 3} \\
RVV & \texttt{001C} & \texttt{006585B3} & \texttt{add a1, a1, t1} \\
RVV & \texttt{0020} & \texttt{40550533} & \texttt{sub a0, a0, t0} \\
RVV & \texttt{0024} & \texttt{FE0512E3} & \texttt{bnez a0, loop} \\
RVV & \texttt{0028} & \texttt{00008067} & \texttt{end: ret} \\
\midrule
Intel & \texttt{0000} & \texttt{C5FB120E} & \texttt{scal64: vmovddup xmm1, qword ptr [rsi]} \\
Intel & \texttt{0004} & \texttt{C5F91007} & \texttt{vmovupd xmm0, [rdi]} \\
Intel & \texttt{0008} & \texttt{C5F959C1} & \texttt{vmulpd xmm0, xmm0, xmm1} \\
Intel & \texttt{000C} & \texttt{C5F91107} & \texttt{vmovupd [rdi], xmm0} \\
Intel & \texttt{0010} & \texttt{4883C710} & \texttt{add rdi, 16} \\
Intel & \texttt{0014} & \texttt{4883EA02} & \texttt{sub rdx, 2} \\
Intel & \texttt{0018} & \texttt{75E6} & \texttt{jnz scal64} \\
\midrule
ARM & \texttt{0000} & \texttt{4D40CC41} & \texttt{scal64: ld1r \{v1.2d\}, [x2]} \\
ARM & \texttt{0004} & \texttt{4C407C20} & \texttt{ld1 \{v0.2d\}, [x1]} \\
ARM & \texttt{0008} & \texttt{6E61DC00} & \texttt{fmul v0.2d, v0.2d, v1.2d} \\
ARM & \texttt{000C} & \texttt{4C007C20} & \texttt{st1 \{v0.2d\}, [x1]} \\
ARM & \texttt{0010} & \texttt{91004021} & \texttt{add x1, x1, \#16} \\
ARM & \texttt{0014} & \texttt{F1000800} & \texttt{subs x0, x0, \#2} \\
ARM & \texttt{0018} & \texttt{54FFFF41} & \texttt{b.ne scal64} \\
\bottomrule
\end{tabular}
\endgroup
\end{table}
\noindent\path{kernels/bx64/scal64.sew64.lmul1.tu.mu.vlen128.xlen64.le.bin}\par

\begin{table}[htbp]
\centering
\caption{Kernel comparison: copy64}
\begingroup
\scriptsize
\setlength{\tabcolsep}{4pt}
\noindent \texttt{a0 = n, a1 = src, a2 = dst}\par
\vspace{0.6ex}
\begin{tabular}{@{}l|l|l|p{7.2cm}@{}}
\toprule
ISA & Address & Hex & Assembly \\
\midrule
RVV & \texttt{0000} & \texttt{02050263} & \texttt{copy64: beqz a0, end} \\
RVV & \texttt{0004} & \texttt{003572D7} & \texttt{loop: vsetvli t0, a0, e64, m1, tu, mu} \\
RVV & \texttt{0008} & \texttt{0205F087} & \texttt{vle64.v v1, (a1)} \\
RVV & \texttt{000C} & \texttt{020670A7} & \texttt{vse64.v v1, (a2)} \\
RVV & \texttt{0010} & \texttt{00329313} & \texttt{slli t1, t0, 3} \\
RVV & \texttt{0014} & \texttt{006585B3} & \texttt{add a1, a1, t1} \\
RVV & \texttt{0018} & \texttt{00660633} & \texttt{add a2, a2, t1} \\
RVV & \texttt{001C} & \texttt{40550533} & \texttt{sub a0, a0, t0} \\
RVV & \texttt{0020} & \texttt{FE0512E3} & \texttt{bnez a0, loop} \\
RVV & \texttt{0024} & \texttt{00008067} & \texttt{end: ret} \\
\midrule
Intel & \texttt{0000} & \texttt{C5F91006} & \texttt{copy64: vmovupd xmm0, [rsi]} \\
Intel & \texttt{0004} & \texttt{C5F91107} & \texttt{vmovupd [rdi], xmm0} \\
Intel & \texttt{0008} & \texttt{4883C610} & \texttt{add rsi, 16} \\
Intel & \texttt{000C} & \texttt{4883C710} & \texttt{add rdi, 16} \\
Intel & \texttt{0010} & \texttt{4883EA02} & \texttt{sub rdx, 2} \\
Intel & \texttt{0014} & \texttt{75EA} & \texttt{jnz copy64} \\
\midrule
ARM & \texttt{0000} & \texttt{4C407C20} & \texttt{copy64: ld1 \{v0.2d\}, [x1]} \\
ARM & \texttt{0004} & \texttt{4C007C40} & \texttt{st1 \{v0.2d\}, [x2]} \\
ARM & \texttt{0008} & \texttt{91004021} & \texttt{add x1, x1, \#16} \\
ARM & \texttt{000C} & \texttt{91004042} & \texttt{add x2, x2, \#16} \\
ARM & \texttt{0010} & \texttt{F1000800} & \texttt{subs x0, x0, \#2} \\
ARM & \texttt{0014} & \texttt{54FFFF61} & \texttt{b.ne copy64} \\
\bottomrule
\end{tabular}
\endgroup
\end{table}
\noindent\path{kernels/bx64/copy64.sew64.lmul1.tu.mu.vlen128.xlen64.le.bin}\par

\begin{table}[htbp]
\centering
\caption{Kernel comparison: axpy64}
\begingroup
\scriptsize
\setlength{\tabcolsep}{4pt}
\noindent \texttt{a0 = n, a1 = x, a2 = y, a3 = alpha. y += alpha*x}\par
\vspace{0.6ex}
\begin{tabular}{@{}l|l|l|p{7.2cm}@{}}
\toprule
ISA & Address & Hex & Assembly \\
\midrule
RVV & \texttt{0000} & \texttt{0006B007} & \texttt{axpy64: fld ft0, 0(a3)} \\
RVV & \texttt{0004} & \texttt{02050663} & \texttt{beqz a0, end} \\
RVV & \texttt{0008} & \texttt{003572D7} & \texttt{loop: vsetvli t0, a0, e64, m1, tu, mu} \\
RVV & \texttt{000C} & \texttt{0205F087} & \texttt{vle64.v v1, (a1)} \\
RVV & \texttt{0010} & \texttt{02067107} & \texttt{vle64.v v2, (a2)} \\
RVV & \texttt{0014} & \texttt{B2105157} & \texttt{vfmacc.vf v2, ft0, v1} \\
RVV & \texttt{0018} & \texttt{02067127} & \texttt{vse64.v v2, (a2)} \\
RVV & \texttt{001C} & \texttt{00329313} & \texttt{slli t1, t0, 3} \\
RVV & \texttt{0020} & \texttt{006585B3} & \texttt{add a1, a1, t1} \\
RVV & \texttt{0024} & \texttt{00660633} & \texttt{add a2, a2, t1} \\
RVV & \texttt{0028} & \texttt{40550533} & \texttt{sub a0, a0, t0} \\
RVV & \texttt{002C} & \texttt{FC051EE3} & \texttt{bnez a0, loop} \\
RVV & \texttt{0030} & \texttt{00008067} & \texttt{end: ret} \\
\midrule
Intel & \texttt{0000} & \texttt{C5FB1211} & \texttt{axpy64: vmovddup xmm2, qword ptr [rcx]} \\
Intel & \texttt{0004} & \texttt{C5F91006} & \texttt{vmovupd xmm0, [rsi]} \\
Intel & \texttt{0008} & \texttt{C5F9100F} & \texttt{vmovupd xmm1, [rdi]} \\
Intel & \texttt{000C} & \texttt{C4E2F9B8CA} & \texttt{vfmadd231pd xmm1, xmm0, xmm2} \\
Intel & \texttt{0011} & \texttt{C5F9110F} & \texttt{vmovupd [rdi], xmm1} \\
Intel & \texttt{0015} & \texttt{4883C610} & \texttt{add rsi, 16} \\
Intel & \texttt{0019} & \texttt{4883C710} & \texttt{add rdi, 16} \\
Intel & \texttt{001D} & \texttt{4883EA02} & \texttt{sub rdx, 2} \\
Intel & \texttt{0021} & \texttt{75DD} & \texttt{jnz axpy64} \\
\midrule
ARM & \texttt{0000} & \texttt{4D40CC62} & \texttt{axpy64: ld1r \{v2.2d\}, [x3]} \\
ARM & \texttt{0004} & \texttt{4C407C20} & \texttt{ld1 \{v0.2d\}, [x1]} \\
ARM & \texttt{0008} & \texttt{4C407C41} & \texttt{ld1 \{v1.2d\}, [x2]} \\
ARM & \texttt{000C} & \texttt{4E62CC01} & \texttt{fmla v1.2d, v0.2d, v2.2d} \\
ARM & \texttt{0010} & \texttt{4C007C41} & \texttt{st1 \{v1.2d\}, [x2]} \\
ARM & \texttt{0014} & \texttt{91004021} & \texttt{add x1, x1, \#16} \\
ARM & \texttt{0018} & \texttt{91004042} & \texttt{add x2, x2, \#16} \\
ARM & \texttt{001C} & \texttt{F1000800} & \texttt{subs x0, x0, \#2} \\
ARM & \texttt{0020} & \texttt{54FFFF01} & \texttt{b.ne axpy64} \\
\bottomrule
\end{tabular}
\endgroup
\end{table}
\noindent\path{kernels/bx64/axpy64.sew64.lmul1.tu.mu.vlen128.xlen64.le.bin}\par

\begin{table}[htbp]
\centering
\caption{Kernel comparison: dot64}
\begingroup
\scriptsize
\setlength{\tabcolsep}{4pt}
\noindent \texttt{a0 = n, a1 = x, a2 = y, a3 = sum out}\par
\vspace{0.6ex}
\begin{tabular}{@{}l|l|l|p{7.2cm}@{}}
\toprule
ISA & Address & Hex & Assembly \\
\midrule
RVV & \texttt{0000} & \texttt{5E003257} & \texttt{dot64: vmv.v.i v4, 0} \\
RVV & \texttt{0004} & \texttt{02050663} & \texttt{beqz a0, end} \\
RVV & \texttt{0008} & \texttt{003572D7} & \texttt{loop: vsetvli t0, a0, e64, m1, tu, mu} \\
RVV & \texttt{000C} & \texttt{0205F087} & \texttt{vle64.v v1, (a1)} \\
RVV & \texttt{0010} & \texttt{02067107} & \texttt{vle64.v v2, (a2)} \\
RVV & \texttt{0014} & \texttt{921110D7} & \texttt{vfmul.vv v1, v1, v2} \\
RVV & \texttt{0018} & \texttt{06121257} & \texttt{vfredusum.vsv4, v1, v4} \\
RVV & \texttt{001C} & \texttt{00329313} & \texttt{slli t1, t0, 3} \\
RVV & \texttt{0020} & \texttt{006585B3} & \texttt{add a1, a1, t1} \\
RVV & \texttt{0024} & \texttt{00660633} & \texttt{add a2, a2, t1} \\
RVV & \texttt{0028} & \texttt{40550533} & \texttt{sub a0, a0, t0} \\
RVV & \texttt{002C} & \texttt{FC051EE3} & \texttt{bnez a0, loop} \\
RVV & \texttt{0030} & \texttt{42401057} & \texttt{end: vfmv.f.s ft0, v4} \\
RVV & \texttt{0034} & \texttt{0006B027} & \texttt{fsd ft0, 0(a3)} \\
RVV & \texttt{0038} & \texttt{00008067} & \texttt{ret} \\
\midrule
Intel & \texttt{0000} & \texttt{C5F91006} & \texttt{dot64: vmovupd xmm0, [rsi]} \\
Intel & \texttt{0004} & \texttt{C5F9100A} & \texttt{vmovupd xmm1, [rdx]} \\
Intel & \texttt{0008} & \texttt{C5F959C1} & \texttt{vmulpd xmm0, xmm0, xmm1} \\
Intel & \texttt{000C} & \texttt{660F7CC0} & \texttt{haddpd xmm0, xmm0} \\
\midrule
ARM & \texttt{0000} & \texttt{4C407C20} & \texttt{dot64: ld1 \{v0.2d\}, [x1]} \\
ARM & \texttt{0004} & \texttt{4C407C41} & \texttt{ld1 \{v1.2d\}, [x2]} \\
ARM & \texttt{0008} & \texttt{6E61DC00} & \texttt{fmul v0.2d, v0.2d, v1.2d} \\
ARM & \texttt{000C} & \texttt{6E60D400} & \texttt{faddp v0.2d, v0.2d, v0.2d} \\
\bottomrule
\end{tabular}
\endgroup
\end{table}
\noindent\path{kernels/bx64/dot64.sew64.lmul1.tu.mu.vlen128.xlen64.le.bin}\par

\begin{table}[htbp]
\centering
\caption{Kernel comparison: nrm264}
\begingroup
\scriptsize
\setlength{\tabcolsep}{4pt}
\noindent \texttt{a0 = n, a1 = x, a2 = norm out. sqrt of sum of squares}\par
\vspace{0.6ex}
\begin{tabular}{@{}l|l|l|p{7.2cm}@{}}
\toprule
ISA & Address & Hex & Assembly \\
\midrule
RVV & \texttt{0000} & \texttt{5E003257} & \texttt{nrm264: vmv.v.i v4, 0} \\
RVV & \texttt{0004} & \texttt{02050263} & \texttt{beqz a0, end} \\
RVV & \texttt{0008} & \texttt{003572D7} & \texttt{loop: vsetvli t0, a0, e64, m1, tu, mu} \\
RVV & \texttt{000C} & \texttt{0205F087} & \texttt{vle64.v v1, (a1)} \\
RVV & \texttt{0010} & \texttt{921090D7} & \texttt{vfmul.vv v1, v1, v1} \\
RVV & \texttt{0014} & \texttt{06121257} & \texttt{vfredusum.vsv4, v1, v4} \\
RVV & \texttt{0018} & \texttt{00329313} & \texttt{slli t1, t0, 3} \\
RVV & \texttt{001C} & \texttt{006585B3} & \texttt{add a1, a1, t1} \\
RVV & \texttt{0020} & \texttt{40550533} & \texttt{sub a0, a0, t0} \\
RVV & \texttt{0024} & \texttt{FE0512E3} & \texttt{bnez a0, loop} \\
RVV & \texttt{0028} & \texttt{42401057} & \texttt{end: vfmv.f.s ft0, v4} \\
RVV & \texttt{002C} & \texttt{5A007053} & \texttt{fsqrt.d ft0, ft0} \\
RVV & \texttt{0030} & \texttt{00063027} & \texttt{fsd ft0, 0(a2)} \\
RVV & \texttt{0034} & \texttt{00008067} & \texttt{ret} \\
\midrule
Intel & \texttt{0000} & \texttt{C5F91006} & \texttt{nrm264: vmovupd xmm0, [rsi]} \\
Intel & \texttt{0004} & \texttt{C5F959C0} & \texttt{vmulpd xmm0, xmm0, xmm0} \\
Intel & \texttt{0008} & \texttt{660F7CC0} & \texttt{haddpd xmm0, xmm0} \\
Intel & \texttt{000C} & \texttt{F20F51C0} & \texttt{sqrtsd xmm0, xmm0} \\
\midrule
ARM & \texttt{0000} & \texttt{4C407C20} & \texttt{nrm264: ld1 \{v0.2d\}, [x1]} \\
ARM & \texttt{0004} & \texttt{6E60DC00} & \texttt{fmul v0.2d, v0.2d, v0.2d} \\
ARM & \texttt{0008} & \texttt{6E60D400} & \texttt{faddp v0.2d, v0.2d, v0.2d} \\
ARM & \texttt{000C} & \texttt{1E61C000} & \texttt{fsqrt d0, d0} \\
\bottomrule
\end{tabular}
\endgroup
\end{table}
\noindent\path{kernels/bx64/nrm264.sew64.lmul1.tu.mu.vlen128.xlen64.le.bin}\par

\begin{table}[htbp]
\centering
\caption{Kernel comparison: asum64}
\begingroup
\scriptsize
\setlength{\tabcolsep}{4pt}
\noindent \texttt{a0 = n, a1 = x, a2 = sum of absolute values}\par
\vspace{0.6ex}
\begin{tabular}{@{}l|l|l|p{7.2cm}@{}}
\toprule
ISA & Address & Hex & Assembly \\
\midrule
RVV & \texttt{0000} & \texttt{5E003257} & \texttt{asum64: vmv.v.i v4, 0} \\
RVV & \texttt{0004} & \texttt{02050263} & \texttt{beqz a0, end} \\
RVV & \texttt{0008} & \texttt{003572D7} & \texttt{loop: vsetvli t0, a0, e64, m1, tu, mu} \\
RVV & \texttt{000C} & \texttt{0205F087} & \texttt{vle64.v v1, (a1)} \\
RVV & \texttt{0010} & \texttt{2A1090D7} & \texttt{vfsgnjx.vv v1, v1, v1} \\
RVV & \texttt{0014} & \texttt{06121257} & \texttt{vfredusum.vsv4, v1, v4} \\
RVV & \texttt{0018} & \texttt{00329313} & \texttt{slli t1, t0, 3} \\
RVV & \texttt{001C} & \texttt{006585B3} & \texttt{add a1, a1, t1} \\
RVV & \texttt{0020} & \texttt{40550533} & \texttt{sub a0, a0, t0} \\
RVV & \texttt{0024} & \texttt{FE0512E3} & \texttt{bnez a0, loop} \\
RVV & \texttt{0028} & \texttt{42401057} & \texttt{end: vfmv.f.s ft0, v4} \\
RVV & \texttt{002C} & \texttt{00063027} & \texttt{fsd ft0, 0(a2)} \\
RVV & \texttt{0030} & \texttt{00008067} & \texttt{ret} \\
\midrule
Intel & \texttt{0000} & \texttt{C5F91006} & \texttt{asum64: vmovupd xmm0, [rsi]} \\
Intel & \texttt{0004} & \texttt{C5F954C7} & \texttt{vandpd xmm0, xmm0, xmm7} \\
Intel & \texttt{0008} & \texttt{660F7CC0} & \texttt{haddpd xmm0, xmm0} \\
\midrule
ARM & \texttt{0000} & \texttt{4C407C20} & \texttt{asum64: ld1 \{v0.2d\}, [x1]} \\
ARM & \texttt{0004} & \texttt{4EE0F800} & \texttt{fabs v0.2d, v0.2d} \\
ARM & \texttt{0008} & \texttt{6E60D400} & \texttt{faddp v0.2d, v0.2d, v0.2d} \\
\bottomrule
\end{tabular}
\endgroup
\end{table}
\noindent\path{kernels/bx64/asum64.sew64.lmul1.tu.mu.vlen128.xlen64.le.bin}\par

\begin{table}[htbp]
\centering
\caption{Kernel comparison: swap64}
\begingroup
\scriptsize
\setlength{\tabcolsep}{4pt}
\noindent \texttt{a0 = n, a1 = x, a2 = y}\par
\vspace{0.6ex}
\begin{tabular}{@{}l|l|l|p{7.2cm}@{}}
\toprule
ISA & Address & Hex & Assembly \\
\midrule
RVV & \texttt{0000} & \texttt{02050663} & \texttt{swap64: beqz a0, end} \\
RVV & \texttt{0004} & \texttt{003572D7} & \texttt{loop: vsetvli t0, a0, e64, m1, tu, mu} \\
RVV & \texttt{0008} & \texttt{0205F087} & \texttt{vle64.v v1, (a1)} \\
RVV & \texttt{000C} & \texttt{02067107} & \texttt{vle64.v v2, (a2)} \\
RVV & \texttt{0010} & \texttt{0205F127} & \texttt{vse64.v v2, (a1)} \\
RVV & \texttt{0014} & \texttt{020670A7} & \texttt{vse64.v v1, (a2)} \\
RVV & \texttt{0018} & \texttt{00329313} & \texttt{slli t1, t0, 3} \\
RVV & \texttt{001C} & \texttt{006585B3} & \texttt{add a1, a1, t1} \\
RVV & \texttt{0020} & \texttt{00660633} & \texttt{add a2, a2, t1} \\
RVV & \texttt{0024} & \texttt{40550533} & \texttt{sub a0, a0, t0} \\
RVV & \texttt{0028} & \texttt{FC051EE3} & \texttt{bnez a0, loop} \\
RVV & \texttt{002C} & \texttt{00008067} & \texttt{end: ret} \\
\midrule
Intel & \texttt{0000} & \texttt{C5F91006} & \texttt{swap64: vmovupd xmm0, [rsi]} \\
Intel & \texttt{0004} & \texttt{C5F9100F} & \texttt{vmovupd xmm1, [rdi]} \\
Intel & \texttt{0008} & \texttt{C5F9110E} & \texttt{vmovupd [rsi], xmm1} \\
Intel & \texttt{000C} & \texttt{C5F91107} & \texttt{vmovupd [rdi], xmm0} \\
Intel & \texttt{0010} & \texttt{4883C610} & \texttt{add rsi, 16} \\
Intel & \texttt{0014} & \texttt{4883C710} & \texttt{add rdi, 16} \\
Intel & \texttt{0018} & \texttt{4883EA02} & \texttt{sub rdx, 2} \\
Intel & \texttt{001C} & \texttt{75E2} & \texttt{jnz swap64} \\
\midrule
ARM & \texttt{0000} & \texttt{4C407C20} & \texttt{swap64: ld1 \{v0.2d\}, [x1]} \\
ARM & \texttt{0004} & \texttt{4C407C41} & \texttt{ld1 \{v1.2d\}, [x2]} \\
ARM & \texttt{0008} & \texttt{4C007C21} & \texttt{st1 \{v1.2d\}, [x1]} \\
ARM & \texttt{000C} & \texttt{4C007C40} & \texttt{st1 \{v0.2d\}, [x2]} \\
ARM & \texttt{0010} & \texttt{91004021} & \texttt{add x1, x1, \#16} \\
ARM & \texttt{0014} & \texttt{91004042} & \texttt{add x2, x2, \#16} \\
ARM & \texttt{0018} & \texttt{F1000800} & \texttt{subs x0, x0, \#2} \\
ARM & \texttt{001C} & \texttt{54FFFF21} & \texttt{b.ne swap64} \\
\bottomrule
\end{tabular}
\endgroup
\end{table}
\noindent\path{kernels/bx64/swap64.sew64.lmul1.tu.mu.vlen128.xlen64.le.bin}\par

\begin{table}[htbp]
\centering
\caption{Kernel comparison: rot64}
\begingroup
\scriptsize
\setlength{\tabcolsep}{4pt}
\noindent \texttt{a0 = n, a1 = x, a2 = y, a3 = (c, s)}\par
\vspace{0.6ex}
\begin{tabular}{@{}l|l|l|p{7.2cm}@{}}
\toprule
ISA & Address & Hex & Assembly \\
\midrule
RVV & \texttt{0000} & \texttt{0006B007} & \texttt{rot64: fld ft0, 0(a3)} \\
RVV & \texttt{0004} & \texttt{0086B087} & \texttt{fld ft1, 8(a3)} \\
RVV & \texttt{0008} & \texttt{04050263} & \texttt{beqz a0, end} \\
RVV & \texttt{000C} & \texttt{003572D7} & \texttt{loop: vsetvli t0, a0, e64, m1, tu, mu} \\
RVV & \texttt{0010} & \texttt{0205F087} & \texttt{vle64.v v1, (a1)} \\
RVV & \texttt{0014} & \texttt{02067107} & \texttt{vle64.v v2, (a2)} \\
RVV & \texttt{0018} & \texttt{921051D7} & \texttt{vfmul.vf v3, v1, ft0} \\
RVV & \texttt{001C} & \texttt{9220D257} & \texttt{vfmul.vf v4, v2, ft1} \\
RVV & \texttt{0020} & \texttt{023212D7} & \texttt{vfadd.vv v5, v3, v4} \\
RVV & \texttt{0024} & \texttt{922051D7} & \texttt{vfmul.vf v3, v2, ft0} \\
RVV & \texttt{0028} & \texttt{9210D257} & \texttt{vfmul.vf v4, v1, ft1} \\
RVV & \texttt{002C} & \texttt{0A321357} & \texttt{vfsub.vv v6, v3, v4} \\
RVV & \texttt{0030} & \texttt{0205F2A7} & \texttt{vse64.v v5, (a1)} \\
RVV & \texttt{0034} & \texttt{02067327} & \texttt{vse64.v v6, (a2)} \\
RVV & \texttt{0038} & \texttt{00329313} & \texttt{slli t1, t0, 3} \\
RVV & \texttt{003C} & \texttt{006585B3} & \texttt{add a1, a1, t1} \\
RVV & \texttt{0040} & \texttt{00660633} & \texttt{add a2, a2, t1} \\
RVV & \texttt{0044} & \texttt{40550533} & \texttt{sub a0, a0, t0} \\
RVV & \texttt{0048} & \texttt{FC0512E3} & \texttt{bnez a0, loop} \\
RVV & \texttt{004C} & \texttt{00008067} & \texttt{end: ret} \\
\midrule
Intel & \texttt{0000} & \texttt{C5FB1211} & \texttt{rot64: vmovddup xmm2, qword ptr [rcx]} \\
Intel & \texttt{0004} & \texttt{C5FB125908} & \texttt{vmovddup xmm3, qword ptr [rcx+8]} \\
Intel & \texttt{0009} & \texttt{C5F91006} & \texttt{vmovupd xmm0, [rsi]} \\
Intel & \texttt{000D} & \texttt{C5F9100F} & \texttt{vmovupd xmm1, [rdi]} \\
Intel & \texttt{0011} & \texttt{C5F959E2} & \texttt{vmulpd xmm4, xmm0, xmm2} \\
Intel & \texttt{0015} & \texttt{C4E2F1B8E3} & \texttt{vfmadd231pd xmm4, xmm1, xmm3} \\
Intel & \texttt{001A} & \texttt{C5F159EA} & \texttt{vmulpd xmm5, xmm1, xmm2} \\
Intel & \texttt{001E} & \texttt{C4E2F9BCEB} & \texttt{vfnmadd231pd xmm5, xmm0, xmm3} \\
Intel & \texttt{0023} & \texttt{C5F91126} & \texttt{vmovupd [rsi], xmm4} \\
Intel & \texttt{0027} & \texttt{C5F9112F} & \texttt{vmovupd [rdi], xmm5} \\
\midrule
ARM & \texttt{0000} & \texttt{4D40CC62} & \texttt{rot64: ld1r \{v2.2d\}, [x3]} \\
ARM & \texttt{0004} & \texttt{91002063} & \texttt{add x3, x3, \#8} \\
ARM & \texttt{0008} & \texttt{4D40CC63} & \texttt{ld1r \{v3.2d\}, [x3]} \\
ARM & \texttt{000C} & \texttt{4C407C20} & \texttt{ld1 \{v0.2d\}, [x1]} \\
ARM & \texttt{0010} & \texttt{4C407C41} & \texttt{ld1 \{v1.2d\}, [x2]} \\
ARM & \texttt{0014} & \texttt{6E62DC04} & \texttt{fmul v4.2d, v0.2d, v2.2d} \\
ARM & \texttt{0018} & \texttt{4E63CC24} & \texttt{fmla v4.2d, v1.2d, v3.2d} \\
ARM & \texttt{001C} & \texttt{6E62DC25} & \texttt{fmul v5.2d, v1.2d, v2.2d} \\
ARM & \texttt{0020} & \texttt{4EE3CC05} & \texttt{fmls v5.2d, v0.2d, v3.2d} \\
ARM & \texttt{0024} & \texttt{4C007C24} & \texttt{st1 \{v4.2d\}, [x1]} \\
ARM & \texttt{0028} & \texttt{4C007C45} & \texttt{st1 \{v5.2d\}, [x2]} \\
\bottomrule
\end{tabular}
\endgroup
\end{table}
\noindent\path{kernels/bx64/rot64.sew64.lmul1.tu.mu.vlen128.xlen64.le.bin}\par

\begin{table}[htbp]
\centering
\caption{Kernel comparison: gemv64}
\begingroup
\scriptsize
\setlength{\tabcolsep}{4pt}
\noindent \texttt{a0 = columns, a1 = A, a2 = x, a3 = y. rows = VLMAX, lda = VLMAX}\par
\vspace{0.6ex}
\begin{tabular}{@{}l|l|l|p{7.2cm}@{}}
\toprule
ISA & Address & Hex & Assembly \\
\midrule
RVV & \texttt{0000} & \texttt{02050863} & \texttt{gemv64: beqz a0, end} \\
RVV & \texttt{0004} & \texttt{003072D7} & \texttt{vsetvli t0, zero, e64, m1, tu, mu} \\
RVV & \texttt{0008} & \texttt{00329313} & \texttt{slli t1, t0, 3} \\
RVV & \texttt{000C} & \texttt{0206F407} & \texttt{vle64.v v8, (a3)} \\
RVV & \texttt{0010} & \texttt{0205F087} & \texttt{loop: vle64.v v1, (a1)} \\
RVV & \texttt{0014} & \texttt{00063007} & \texttt{fld ft0, 0(a2)} \\
RVV & \texttt{0018} & \texttt{B2105457} & \texttt{vfmacc.vf v8, ft0, v1} \\
RVV & \texttt{001C} & \texttt{006585B3} & \texttt{add a1, a1, t1} \\
RVV & \texttt{0020} & \texttt{00860613} & \texttt{addi a2, a2, 8} \\
RVV & \texttt{0024} & \texttt{FFF50513} & \texttt{addi a0, a0, -1} \\
RVV & \texttt{0028} & \texttt{FE0514E3} & \texttt{bnez a0, loop} \\
RVV & \texttt{002C} & \texttt{0206F427} & \texttt{vse64.v v8, (a3)} \\
RVV & \texttt{0030} & \texttt{00008067} & \texttt{end: ret} \\
\midrule
Intel & \texttt{0000} & \texttt{C5F91006} & \texttt{gemv64: vmovupd xmm0, [rsi]} \\
Intel & \texttt{0004} & \texttt{C5FB120A} & \texttt{vmovddup xmm1, qword ptr [rdx]} \\
Intel & \texttt{0008} & \texttt{C4E2F9B8E1} & \texttt{vfmadd231pd xmm4, xmm0, xmm1} \\
Intel & \texttt{000D} & \texttt{4883C610} & \texttt{add rsi, 16} \\
Intel & \texttt{0011} & \texttt{4883C208} & \texttt{add rdx, 8} \\
Intel & \texttt{0015} & \texttt{4883E901} & \texttt{sub rcx, 1} \\
Intel & \texttt{0019} & \texttt{75E5} & \texttt{jnz gemv64} \\
\midrule
ARM & \texttt{0000} & \texttt{4CDF7C20} & \texttt{gemv64: ld1 \{v0.2d\}, [x1], \#16} \\
ARM & \texttt{0004} & \texttt{4DDFCC41} & \texttt{ld1r \{v1.2d\}, [x2], \#8} \\
ARM & \texttt{0008} & \texttt{4E61CC04} & \texttt{fmla v4.2d, v0.2d, v1.2d} \\
ARM & \texttt{000C} & \texttt{F1000400} & \texttt{subs x0, x0, \#1} \\
ARM & \texttt{0010} & \texttt{54FFFF81} & \texttt{b.ne gemv64} \\
\bottomrule
\end{tabular}
\endgroup
\end{table}
\noindent\path{kernels/bx64/gemv64.sew64.lmul1.tu.mu.vlen128.xlen64.le.bin}\par

\begin{table}[htbp]
\centering
\caption{Kernel comparison: ger64}
\begingroup
\scriptsize
\setlength{\tabcolsep}{4pt}
\noindent \texttt{a0 = columns, a1 = x, a2 = y, a3 = A. A += x y\textasciicircum{}T, lda = VLMAX}\par
\vspace{0.6ex}
\begin{tabular}{@{}l|l|l|p{7.2cm}@{}}
\toprule
ISA & Address & Hex & Assembly \\
\midrule
RVV & \texttt{0000} & \texttt{02050863} & \texttt{ger64: beqz a0, end} \\
RVV & \texttt{0004} & \texttt{003072D7} & \texttt{vsetvli t0, zero, e64, m1, tu, mu} \\
RVV & \texttt{0008} & \texttt{00329313} & \texttt{slli t1, t0, 3} \\
RVV & \texttt{000C} & \texttt{0205F087} & \texttt{vle64.v v1, (a1)} \\
RVV & \texttt{0010} & \texttt{0206F407} & \texttt{loop: vle64.v v8, (a3)} \\
RVV & \texttt{0014} & \texttt{00063007} & \texttt{fld ft0, 0(a2)} \\
RVV & \texttt{0018} & \texttt{B2105457} & \texttt{vfmacc.vf v8, ft0, v1} \\
RVV & \texttt{001C} & \texttt{0206F427} & \texttt{vse64.v v8, (a3)} \\
RVV & \texttt{0020} & \texttt{006686B3} & \texttt{add a3, a3, t1} \\
RVV & \texttt{0024} & \texttt{00860613} & \texttt{addi a2, a2, 8} \\
RVV & \texttt{0028} & \texttt{FFF50513} & \texttt{addi a0, a0, -1} \\
RVV & \texttt{002C} & \texttt{FE0512E3} & \texttt{bnez a0, loop} \\
RVV & \texttt{0030} & \texttt{00008067} & \texttt{end: ret} \\
\midrule
Intel & \texttt{0000} & \texttt{C5F9100E} & \texttt{ger64: vmovupd xmm1, [rsi]} \\
Intel & \texttt{0004} & \texttt{C5F91001} & \texttt{vmovupd xmm0, [rcx]} \\
Intel & \texttt{0008} & \texttt{C5FB1212} & \texttt{vmovddup xmm2, qword ptr [rdx]} \\
Intel & \texttt{000C} & \texttt{C4E2F1B8C2} & \texttt{vfmadd231pd xmm0, xmm1, xmm2} \\
Intel & \texttt{0011} & \texttt{C5F91101} & \texttt{vmovupd [rcx], xmm0} \\
Intel & \texttt{0015} & \texttt{4883C110} & \texttt{add rcx, 16} \\
Intel & \texttt{0019} & \texttt{4883C208} & \texttt{add rdx, 8} \\
Intel & \texttt{001D} & \texttt{4883E801} & \texttt{sub rax, 1} \\
Intel & \texttt{0021} & \texttt{75DD} & \texttt{jnz ger64} \\
\midrule
ARM & \texttt{0000} & \texttt{4C407C21} & \texttt{ger64: ld1 \{v1.2d\}, [x1]} \\
ARM & \texttt{0004} & \texttt{4C407C60} & \texttt{ld1 \{v0.2d\}, [x3]} \\
ARM & \texttt{0008} & \texttt{4DDFCC42} & \texttt{ld1r \{v2.2d\}, [x2], \#8} \\
ARM & \texttt{000C} & \texttt{4E62CC20} & \texttt{fmla v0.2d, v1.2d, v2.2d} \\
ARM & \texttt{0010} & \texttt{4C9F7C60} & \texttt{st1 \{v0.2d\}, [x3], \#16} \\
ARM & \texttt{0014} & \texttt{F1000400} & \texttt{subs x0, x0, \#1} \\
ARM & \texttt{0018} & \texttt{54FFFF41} & \texttt{b.ne ger64} \\
\bottomrule
\end{tabular}
\endgroup
\end{table}
\noindent\path{kernels/bx64/ger64.sew64.lmul1.tu.mu.vlen128.xlen64.le.bin}\par

\begin{table}[htbp]
\centering
\caption{Kernel comparison: iamax64}
\begingroup
\scriptsize
\setlength{\tabcolsep}{4pt}
\noindent \texttt{a0 = n (<= VLMAX), a1 = x, a2 = index out. First index of max |x|}\par
\vspace{0.6ex}
\begin{tabular}{@{}l|l|l|p{7.2cm}@{}}
\toprule
ISA & Address & Hex & Assembly \\
\midrule
RVV & \texttt{0000} & \texttt{02050263} & \texttt{iamax64: beqz a0, end} \\
RVV & \texttt{0004} & \texttt{003572D7} & \texttt{vsetvli t0, a0, e64, m1, tu, mu} \\
RVV & \texttt{0008} & \texttt{0205F087} & \texttt{vle64.v v1, (a1)} \\
RVV & \texttt{000C} & \texttt{2A1090D7} & \texttt{vfsgnjx.vv v1, v1, v1} \\
RVV & \texttt{0010} & \texttt{1E109157} & \texttt{vfredmax.vs v2, v1, v1} \\
RVV & \texttt{0014} & \texttt{42201057} & \texttt{vfmv.f.s ft0, v2} \\
RVV & \texttt{0018} & \texttt{621051D7} & \texttt{vmfeq.vf v3, v1, ft0} \\
RVV & \texttt{001C} & \texttt{4238A2D7} & \texttt{vfirst.m t0, v3} \\
RVV & \texttt{0020} & \texttt{00562023} & \texttt{sw t0, 0(a2)} \\
RVV & \texttt{0024} & \texttt{00008067} & \texttt{end: ret} \\
\midrule
Intel & \texttt{0000} & \texttt{C5F91007} & \texttt{iamax64: vmovupd xmm0, [rdi]} \\
Intel & \texttt{0004} & \texttt{C5F954C7} & \texttt{vandpd xmm0, xmm0, xmm7} \\
Intel & \texttt{0008} & \texttt{0F12C8} & \texttt{movhlps xmm1, xmm0} \\
Intel & \texttt{000B} & \texttt{660F5FC1} & \texttt{maxpd xmm0, xmm1} \\
\midrule
ARM & \texttt{0000} & \texttt{4C407C20} & \texttt{iamax64: ld1 \{v0.2d\}, [x1]} \\
ARM & \texttt{0004} & \texttt{4EE0F800} & \texttt{fabs v0.2d, v0.2d} \\
ARM & \texttt{0008} & \texttt{7E70F801} & \texttt{fmaxp d1, v0.2d} \\
ARM & \texttt{000C} & \texttt{4E080421} & \texttt{dup v1.2d, v1.d[0]} \\
ARM & \texttt{0010} & \texttt{6EE18C02} & \texttt{cmeq v2.2d, v0.2d, v1.2d} \\
\bottomrule
\end{tabular}
\endgroup
\end{table}
\noindent\path{kernels/bx64/iamax64.sew64.lmul1.tu.mu.vlen128.xlen64.le.bin}\par


\section{Harts}

The organisation that justifies a second hart is the one the wide machines are for: one hart runs the operating system and short vector work, the other runs a longer kernel from a local memory. Both are RV64 with the same vector claim. A DMA engine moves blocks into that local memory. The DMA is a bus master programmed by stores, in the same style as the blitter on BX-16. It does not fetch vector instructions.

Local memory is mapped. Vector loads and stores name it with ordinary addresses. A kernel that ran on BX-32 in SEW=32 runs here in SEW=32 without change. A kernel that wants FP64 sets SEW=64 and uses the double-precision opcodes.

\section{Xsvm on this machine}

Leave Xsvm off on the image you use to claim V. A standard toolchain will not emit it, and a custom opcode in the middle of an otherwise clean V implementation confuses every test suite you did not write. A second elaboration, with Xsvm on, is how you compare a four-by-four custom instruction with the fused-multiply sequence. Publish both cycle counts. Keep the V-claiming image free of custom geometry.
